\chapter*{Abstract}

Electroencephalography (EEG) is a central tool in clinical and experimental neuroscience, but channel loss caused by contact failures, motion artifacts, or amplifier saturation can compromise downstream analyses. This thesis evaluates methods for interpolating missing EEG channels under a reproducible protocol that includes geometric, statistical, Graph Signal Processing (GSP), and spatio-temporal regularization approaches.

The study uses three public datasets (PhysioNet, BCI Competition IV 2a, and MNE Sample), random and contiguous loss scenarios with severities from 10\% to 40\%, and metrics covering amplitude error, temporal shape, spectral preservation, and computational cost. The evaluation is organized into phases: preliminary method--graph selection, candidate reduction, intermediate optimization, and a final paired comparison between TRSS and MNE-Python.

The results support a conditional conclusion. Fixed TRSS reduces median MAE by 12.4\% relative to MNE-Python, reaches a 72.0\% win rate, and improves NRMSE by 13.9\%. Its advantage increases under severe nearby or peripheral losses, where temporal continuity compensates for missing local spatial information. However, MNE-Python retains advantages in computational cost, operational stability, and global spectral preservation. Thus, TRSS is advisable when reconstruction fidelity under difficult spatial loss is the priority and offline processing is acceptable; MNE-Python remains the robust reference for standard preprocessing and time-constrained scenarios.

Code, configurations, optimization logs, and result artifacts are preserved to ensure full traceability.

\vspace{1em}
\noindent\textbf{Keywords:} EEG, missing channels, interpolation, Graph Signal Processing, temporal regularization, TRSS, MNE-Python, reproducible evaluation.